\section{模型假设}
\begin{table}[H]\centering\caption{模型假设、依据和失效条件}
\begin{tabularx}{\linewidth}{lYY}\toprule
编号 & 假设与使用位置 & 可能失效及检查方式\\\midrule
H1 & 输入图为题面给定的 DAG；节点 id 在图内唯一 & 若依赖成环，方案合法性检查必须失败并保留日志\\
H2 & Tensor 大小由原始 JSON 的 `size` 字段给出，边本身不重复计量 & 读取或单位异常时回查原始文件和配置\\
H3 & L1/UB 容量、DDR 带宽、同步延迟和 Cache 参数以 `config.txt` 与题面为准 & 任何运行改变配置都不能与正式结果合并\\
H4 & 轻量评分只作候选排序，官方核内和多核评估才给出正式目标 & 代理排序与真实 Makespan 不一致时保留反例\\
H5 & 同一规范化计划在相同图、场景、核数和配置下可复用 & 计划哈希或代码指纹变化时重新评估\\
H6 & 问题三的 Cache 按逻辑 tensor id 查询并遵循题面 FIFO 规则 & id 缺失、单张量超过容量或事件顺序异常时标记不适用\\\bottomrule
\end{tabularx}\end{table}

上述假设只用于定义算法和证据边界，不直接推出任何性能结论。尤其是有限候选搜索不能推出全局最优；程序通过测试也不能替代队员对公式、单位、结果和题面要求的人工核验。实验完成后，稳健性检查将优先覆盖候选失败、容量退化、核数变化和共享 Cache 查询顺序变化。
